%*******************************************************
% KI-Verzeichnis
% (Offenlegung der eingesetzten KI-basierten Hilfsmittel)
% [Nicht übersetzen, da Prüfungsleistung an dt. Hochschule]
%*******************************************************
\pdfbookmark[0]{KI-Verzeichnis}{kiverzeichnis}
\chapter*{KI-Verzeichnis}

\noindent Die folgende Tabelle legt die in dieser Arbeit eingesetzten KI-basierten Hilfsmittel sowie Art und Umfang ihrer Nutzung offen. Die Ideen dahinter stammen vom Verfasser, die Hilfsmittel nehmen ihm diese nicht ab. Ihre Vorschläge wurden nur übernommen, wenn sie passten und nach den Ideen des Verfassers liefen. In den meisten Fällen wurden diese dann noch einmal ohne KI angepasst. Die inhaltliche Verantwortung für die gesamte Arbeit liegt beim Verfasser.
\bigskip

\begin{center}
\renewcommand{\arraystretch}{1.35}
\setlength{\tabcolsep}{5pt}
{\small
\begin{tabular}{c >{\raggedright\arraybackslash}p{2.6cm} >{\raggedright\arraybackslash}p{5.6cm} >{\raggedright\arraybackslash}p{3.9cm}}
\toprule
\textbf{Nr.} & \textbf{KI-Werkzeug} & \textbf{Art der Nutzung} & \textbf{Betroffene Abschnitte} \\
\midrule
1 & Claude (Anthropic)
  & Rechtschreib- und Grammatikkorrektur sowie verbesserte Formulierungsvorschläge. Unterstützung bei der Konzeption von Entwürfen für Kapitel.
  & Gesamte Arbeit \\
\addlinespace
2 & Claude (Anthropic)
  & Entwürfe für Abbildungen, Listings, Tabellen und Formeln nach eigenen Vorgaben, anschließend vom Verfasser überarbeitet. Idee und Aussage stammen ebenso vom Verfasser.
  & Abb.~4.x, 5.1, 5.2, 5.5--5.8, 6.1, 6.2, 6.6 \newline Tab.~6.x, C.1, D.x \newline For.~2.2, 4.x \\
\addlinespace
3 & Claude (Anthropic)
  & Unterstützung bei Programmcode durch Vorschläge (u.\,a. Pseudocode), Korrekturen und Optimierungen sowie Erstellung der Kommentare und Aufräumen des Quellcodes.
  & Gekennzeichnete Stellen \newline Kommentare und Aufräumen im gesamten Quellcode der Erweiterung \\
\addlinespace
4 & DeepL (DeepL SE)
  & Übersetzung von Textpassagen aus dem Englischen ins Deutsche.
  & Alle englischsprachigen Quellen (Literaturverzeichnis) \\
\addlinespace
5 & Perplexity (Perplexity AI) / NotebookLM (Google)
  & Quellenrecherche (Perplexity) und Auffinden von Textstellen in den Quellen (NotebookLM, nur vereinzelt) als erster Schritt, anschließend vom Verfasser in der Originalquelle geprüft.
  & Literaturverzeichnis \newline Quellenbelege \\
\bottomrule
\end{tabular}
}
\end{center}
